\section{Delta Model Module}

\subsection{Purpose}

The model module applies a first-order differential (delta) transform to a byte stream
\emph{before} LZSS encoding, and its exact inverse \emph{after} LZSS decoding.

In a smooth grayscale image, adjacent pixels differ by only a few intensity levels.
The raw pixel stream therefore contains values distributed across the full 0--255 range.
After delta coding, the stream consists mostly of small values near zero
(e.g., $-3, 0, +2, 0, +1, \ldots$), which means many bytes repeat the same small value.
LZSS finds longer matches in such a stream, producing a higher compression ratio.

The effect is strongest for images with gradual gradients and least effective for
high-frequency noise images.

\subsection{Files}

\begin{itemize}
  \item \texttt{include/model.hpp} --- function declarations.
  \item \texttt{src/model.cpp} --- implementations (\textasciitilde{}20 lines total).
\end{itemize}

\subsection{forward\_model()}

\begin{lstlisting}[style=cpp]
void forward_model(std::vector<uint8_t>& data);
\end{lstlisting}

\textbf{Effect:} transforms \texttt{data} in-place so that
$\mathtt{data}[i] = \mathtt{data}[i] - \mathtt{data}[i-1]$ (unsigned 8-bit arithmetic)
for all $i > 0$.
\texttt{data[0]} is left unchanged (it is the anchor value).

\textbf{Loop direction --- backward:}

\begin{lstlisting}[style=cpp]
for (size_t i = data.size() - 1; i > 0; --i)
    data[i] = (uint8_t)(data[i] - data[i - 1]);
\end{lstlisting}

The loop iterates \emph{from the end toward the start}.
This is essential: when computing \texttt{data[i]}, we need the \emph{original}
\texttt{data[i-1]}, not the already-overwritten delta value.
Iterating backward ensures \texttt{data[i-1]} has not yet been modified.

A forward loop would require a temporary variable to preserve the previous original value.
The backward loop achieves the same result with no extra storage.

\textbf{Arithmetic:}
The subtraction uses \texttt{uint8\_t} (unsigned 8-bit) arithmetic.
Subtraction wraps modulo 256 on all platforms without undefined behaviour.
Values of the original stream are in 0--255; differences can be anywhere in 0--255
when interpreted as unsigned (e.g., $3 - 200 = 59$ in unsigned 8-bit arithmetic).

\subsection{inverse\_model()}

\begin{lstlisting}[style=cpp]
void inverse_model(std::vector<uint8_t>& data);
\end{lstlisting}

\textbf{Effect:} reconstructs the original values from a delta-coded buffer.
For all $i > 0$: $\mathtt{data}[i] = \mathtt{data}[i] + \mathtt{data}[i-1]$.
\texttt{data[0]} is unchanged.

\textbf{Loop direction --- forward:}

\begin{lstlisting}[style=cpp]
for (size_t i = 1; i < data.size(); ++i)
    data[i] = (uint8_t)(data[i] + data[i - 1]);
\end{lstlisting}

Here the loop must iterate \emph{forward}: to reconstruct \texttt{data[i]}, we need
the already-reconstructed \texttt{data[i-1]} (its original value, now restored).
Using the reconstructed prefix accumulates the cumulative sum correctly
(prefix-sum / running total).

\subsection{Round-Trip Guarantee}

$\texttt{inverse\_model}(\texttt{forward\_model}(\mathbf{x})) = \mathbf{x}$
for all inputs $\mathbf{x}$.

This holds because modular arithmetic satisfies
$(a - b) + b \equiv a \pmod{256}$.
No floating-point rounding or overflow is possible.

\subsection{Integration with the Codec}

\texttt{forward\_model} is called inside \texttt{codec.cpp} after scanning (scan gives the
1-D byte stream) and before LZSS encoding.
\texttt{inverse\_model} is called after LZSS decoding and before unscanning.

In adaptive mode, the model is applied per-block (each block is scanned independently,
then modelled independently).
The first pixel of each block serves as its own anchor; there is no cross-block
delta dependency.
